\section{Other Backends}
\label{sec:other-backends}

Past the EVM and Gambit backends, the rest of the family share the
EventGraph but address different consumers: a constraint solver
(SMT-LIB), a session-typed protocol checker (Scribble), a Bitcoin
Lightning channel (for the two-party fragment), a MAID exporter
(for the Thrones game-theory workbench), a GraphViz visualisation,
and a Coq/Lean deep embedding of the protocol as a dependent
typing problem.  Each is short enough to cover in one section.

\subsection{SMT-LIB}

The SMT backend (\file{backend/smt/Smt.kt}) encodes the EventGraph
as a set of quantifier-free constraints over per-node ``done''
flags and per-field value variables.  A bounded-type field is a
finite-sort variable; a guard is its boolean lowering;
dependencies become implications
\(\mathit{done}_B \Rightarrow \bigwedge_{A \in \mathrm{pred}(B)} \mathit{done}_A\).
Reachability queries are written as additional constraints
asserting that a chosen role gets a payoff at least $x$.

The result is given to Z3 (or any SMT-LIB-2 solver) for
satisfiability.  Two principal uses:

\begin{itemize}\itemsep1pt
\item \emph{Witness searches.}  ``Is there a play that gives role
  $R$ a payoff $\geq c$?''  A satisfying assignment is a concrete
  history.
\item \emph{Coalition queries.}  A DQBF encoding is also generated,
  in which a chosen coalition's variables are existentially
  quantified and the complement's are universally quantified.  The
  resulting formula is true iff the coalition has a strategy that
  achieves a property against any response.  This is a coarse
  analogue of best-response analysis useful for sanity checks on
  small examples.
\end{itemize}

Z3 is fast on the example library; the encoding is plain enough
that it is also useful as a sanity-check on the IR (a guard that
nobody can ever satisfy will show up as a falsifying constraint).

\subsection{Scribble}

Scribble~\citep{HondaYoshidaCarbone2008MPST} is a session-types
implementation that records the legal sequencing of messages in a
multi-party protocol.  The Scribble backend
(\file{backend/scribble/}) projects the EventGraph onto the
public-message subset --- skipping commits, retaining reveals and
public yields, dropping internal randomness --- and emits a global
type whose sequencing is the topological order of the public
frontier.  Recursion in the produced Scribble protocol is empty
(the protocol is acyclic by construction).

The output is useful to anyone who wants a static check that
their off-chain client library is sending messages in the right
order; Scribble's projection toolchain can derive per-role local
types from the global protocol.

\subsection{Bitcoin / Lightning (sketch)}

A proof-of-concept backend (\file{backend/bitcoin/}) targets the
two-party Lightning channel for the subset of Vegas games with
exactly two strategic roles.  Each EventGraph node is mapped to
a state in a Bitcoin script state machine; commitments use
\TT{OP\_HASH160} and \TT{OP\_EQUAL}, with state transitions
encoded as signed updates of the channel's commitment
transaction.  The output is a textual transcript of the
Lightning protocol moves needed to realise the game --- not a
deployable artifact.  Deployment on a real Lightning node, and
the integration work that would entail, is out-of-band.

Its value here is the existence proof: the EventGraph
abstraction lifts to a non-EVM execution layer with completely
different cryptographic primitives, with no change to the IR.
A proper Bitcoin/Lightning backend is on the future-work list,
not a claim of the present work.

\subsection{MAID}

MAID~\citep{koller2003multi} (Multi-Agent Influence Diagrams)
factor a multi-agent decision problem into decision nodes,
chance nodes, utility nodes, and the conditional probability
distributions that wire them together.  Vegas can emit a MAID
suitable for the Thrones game-theory workbench.

The MAID backend has a real expressiveness limit.  In plain
MAIDs, a decision node's action set is a single static domain;
there is no place to say ``the legal action set depends on the
parents' values''.  A Vegas guard like
\TT{where Host.goat != Guest.d} has nowhere to go in such a
MAID: the encoding would silently drop the constraint and present
the Host as being able to pick any door.  The current MAID
emitter therefore refuses to emit for any game whose guards
reference fields written by other nodes; it accepts self-only
guards (a guard that only constrains the value being written) but
does not yet narrow the resulting MAID node's domain.  Both
narrowing self-only guards and supporting context-dependent
domains in some MAID extension are listed in \file{TODO.txt}.

Where MAID does apply --- protocols with trivial or self-only
guards --- the MAID view is more compact than the unfolded
extensive-form tree.  For the rest, Gambit's EFG is the more
expressive target.

\subsection{Coq / Lean: Diamond DAG witness records}
\label{sec:other:diamond}

The Gallina (Coq) and Lean~4 backends emit a deep embedding of
the EventGraph as a dependently-typed family of records.  The
encoding is worth describing directly: every move and every
guard is a type-level proof obligation, and player strategies
become functions on typed prior witnesses.

For an EventGraph with $n$ nodes, the backend emits a chain of
records $W_0, W_1, \dots, W_n$, where $W_i$ is the (typed) state
after the first $i$ nodes have fired.  Each $W_i$ is parameterised
on the entire history $W_0, \dots, W_{i-1}$, with fields:

\begin{itemize}\itemsep1pt
\item one named field per parameter written by node $i$, of the
  parameter's bounded type, with the field name embedding the
  parameter and owning role (e.g.\ \TT{hidden\_car\_Host} for a
  committed value, \TT{d\_Guest} for a public yield);
\item one proof field per guard ---
  \TT{$W_i$\_guard\_domain\_\textit{f}} for the domain restriction
  and \TT{$W_i$\_guard\_logic} for the user-supplied \TT{where}
  clause;
\item for reveals, a proof
  \TT{$W_i$\_guard\_reveal\_\textit{f}} that the revealed value
  equals the value committed in an ancestor witness;
\item nothing else.
\end{itemize}

The backend emits one of two top-level records, named for the
side of the protocol they encode.  \TT{ActionDag} is the
\emph{fair-play structure}: an inhabitant is a complete play of
the protocol together with proofs of every guard, every domain
restriction, and every reveal-matches-commit equation.  Inhabiting
\TT{ActionDag} is therefore the same as exhibiting that a fair
play exists.  \TT{EventDag} is the \emph{trace structure}: an
inhabitant is a realised history, with hidden fields wrapped in a
\TT{Maybe}-style carrier so that traces can be partially observed.
Actor strategies are functions from the prior witnesses
$W_0, \dots, W_{i-1}$ (projected to the actor's view) into the
appropriate field of $W_i$, with proofs --- so the encoding
captures \emph{at the type level} that every move is legal and
every reveal matches its commit.

The Coq and Lean backends each take a policy flag selecting which
records to emit:

\begin{itemize}\itemsep1pt
\item \TT{FAIR\_PLAY} (\TT{lean-fair} / \TT{gallina-fair}) emits
  \TT{ActionDag} only.  Every guard becomes a runtime
  \TT{decide}/\TT{rfl} obligation in the appropriate $W_i$.
\item \TT{INDEPENDENT} (\TT{lean-independent} /
  \TT{gallina-independent}) emits \TT{EventDag} only, with
  hidden fields wrapped in a \TT{Maybe} carrier and an
  \TT{isPresent} discipline on reveals.
\item \TT{MONOTONIC} (\TT{lean-monotone} / \TT{gallina-monotone})
  emits both records, with monotonicity guards
  (\TT{$W_i$\_guard\_have\_$w_j$}) certifying that prior commits
  remain valid for later reveals.
\end{itemize}

Each policy is exercised against the same example set in the
golden-master harness.

The Diamond DAG encoding is independent of the companion
VegasCore Lean library described in \S\ref{sec:vegascore}.  The
two coexist as different approaches to ``the same protocol viewed
through a proof assistant'': Diamond DAG ships a standalone
per-program object whose inhabitation is a witness; VegasCore is
a reusable library of definitions and theorems.  A third Lean
backend, emitting a Vegas program as an instance of VegasCore's
\TT{WFProgram} type, is planned (and listed in
\file{TODO.txt}) to bridge the two.

The encoding is interesting for two reasons.  First, it is
self-contained: the Lean object can be passed to a theorem
prover without depending on a large semantics library; the
encoding \emph{is} the semantics.  Second, because each step is
a function on prior witnesses, the player's choice schedule
appears as a strategy in the strict type-theoretic sense --- an
inhabitant of a function space whose codomain is pinned by
domain and where-clause proofs.  This view is closer to the
IF-logic and compositional-game-theory style of game
semantics~\citep{HintikkaSandu1989,GhaniHedgesWinschel2018CGT}
than to the state-machine view that Gambit's tree presents.

\subsection{GraphViz}

The simplest backend: a textual DOT description of the EventGraph
for debugging and presentation.  Player nodes are boxes; chance
roles are diamonds.  Visibility is encoded in fill colour: pink
for \COMMIT writes, green for \REVEAL writes, light grey for
\PUBLIC writes.  Useful for inspecting the result of the
commit-reveal expansion pass.
